\section{Conclusion}
\label{sec:conclusion}

We studied arbitrary single-point evaluation of degree-\(<N\) univariate
polynomials represented in the Boolean-kernel basis
\[
K_y(X)
=
\prod_{i=0}^{m-1}
\left(X^{2^i}+y_i\right),
\qquad
N=2^m,
\qquad
y\in\{0,1\}^m.
\]

The representation admits the local fold
\[
F_t(a,b)=t(a+b)+b,
\qquad
t=x^{2^i},
\]
which evaluates the polynomial through a complete binary reduction tree.
The tree contains exactly
\[
N-1
\]
folds and therefore requires exactly
\[
N-1
\]
field multiplications and
\[
2(N-1)
\]
field additions, together with
\[
m-1
\]
squarings to generate the level parameters.

On the implementation side, the final Rust evaluator combines
coarse-grained subtree parallelism, specialized Goldilocks reduction,
a fused kernel fold, bounds-check elimination in the hot loop, direct
first-level reduction, and worker-local scratch reuse.

For
\[
N=2^{20},
\]
the final same-run benchmark on the tested Apple M1 system measured
\[
1.0560\ \mathrm{ms}
\]
for the optimized Boolean-kernel evaluator and
\[
1.4116\ \mathrm{ms}
\]
for the blocked parallel monomial baseline.  This corresponds to a
\[
\boxed{1.337\times}
\]
evaluation-speed advantage, or approximately
\[
\boxed{25.2\%}
\]
lower wall-clock time for the isolated arbitrary-point evaluation
primitive.

A separate kernel-only optimization run reached
\[
892.57\,\mu\mathrm{s},
\]
but we do not combine that number with competitor timings collected in a
different run when computing cross-basis speedups.

The comparison with the roots-of-unity Lagrange representation shows an
even larger difference for arbitrary off-domain evaluation.  At
\(N=2^{20}\), the measured parallel Lagrange evaluator required
\(17.199\) ms.  This does not imply that evaluation form is generally
inferior in proof systems; rather, it emphasizes that representation choice
is workload-dependent.

The central systems question is therefore not whether the Boolean-kernel
basis supports fast single-point evaluation---the experiments establish that
it does---but whether a prover can retain this representation across enough
of its pipeline that conversion, encoding, commitment, and FRI-related
costs do not erase the local advantage.

Future work should measure that representation lifetime directly inside a
complete transparent proof system, including basis conversion, repeated
off-domain queries, commitment, and FRI integration.  Only such an
end-to-end study can determine how much of the isolated evaluator speedup
survives at prover scale.
